\chapter{预测模型、泛化与不确定性}
\label{ch:week07-prediction}

管理决策通常发生在结果出现之前。为了安排库存、维修或客户服务，组织需要根据
当前状态判断未来风险。预测模型近似表达当前信息与未来结果的条件关系。
预测任务应明确预测时点、对象、可用信息、目标结果和评价错误的损失函数。

总体风险衡量目标分布下的预测损失，经验风险以样本平均近似这一期望；逻辑回归
与集成学习是常用的预测方法。样本划分、概率评价、预测区间和运行监测用于
评价模型的泛化及适用范围。预测概率可以进入期望损失计算；改变结果的干预
则需估计处理效应，见第8章。

\section{经营预警与预测任务}

企业管理经常需要提前判断未来状态。例如，哪些客户可能在下月流失，哪些订单可能延期，哪些产品下周需求会上升，哪些设备可能发生故障，哪些供应商可能无法准时履约。这类问题与第~\ref{ch:week06-ab-testing}~章 A/B 测试不同。A/B 测试关注某个行动是否造成结果变化；预测模型关注在当前信息下，未来结果最可能是什么。

设 \(X\) 表示当前可观测的企业状态，\(Y\) 表示未来经营结果。预测问题可以写成
\[
  Y=f(X)+\varepsilon,
\]
在平方损失下，取 \(f(X)=\E[Y\mid X]\)，则残差 \(\varepsilon=Y-f(X)\) 满足 \(\E[\varepsilon\mid X]=0\)。预测模型用样本估计该条件关系，以减小目标分布下的新样本预测误差。

\begin{definition}[监督学习]
监督学习以带标签样本
\[
  \mathcal D=\{(x_i,y_i):i=1,\ldots,n\}
\]
为输入，学习一个映射 \(\hat f:\mathcal X\rightarrow \mathcal Y\)。当 \(Y\) 为连续数值时称为回归问题；当 \(Y\) 为类别或二元结果时称为分类问题。
\end{definition}

客户流失预测就是一个二元分类问题。对客户 \(i\)，令 \(Y_i=1\) 表示在未来观察窗口内流失，\(Y_i=0\) 表示未流失。特征 \(X_i\) 可以包括客户关系时长、月消费、客服工单数、折扣率、使用活跃度和履约延迟次数。模型输出客户在给定信息下的流失概率
\[
  \hat p_i=\widehat{\Prob}(Y_i=1\mid X_i=x_i).
\]
管理者随后还需要决定：当 \(\hat p_i\) 多大时需要人工复核、发放优惠、安排客户经理联系，或者暂不行动。

\section{预测对象与可用信息}
\subsection{预测、解释与因果推断}

预测关联描述给定信息下的结果分布。例如，折扣率与流失风险正相关，可能源于高风险客户更常获得优惠。折扣率可作为既有服务策略下的预测变量；研究提高折扣的效果，则需要比较不同干预下的潜在结果。

\begin{table}[htbp]
\centering
\caption{几类分析任务的边界}
\begin{tabular}{P{0.20\linewidth}P{0.36\linewidth}P{0.32\linewidth}}
\toprule
\rowcolor{BDMTableHead}
任务 & 典型问题 & 主要输出\\
\midrule

描述分析 & 当前发生了什么？ & KPI、分布、趋势和异常。\\
统计推断 & 样本差异是否足以支持总体判断？ & 标准误、置信区间和检验。\\
预测建模 & 哪些对象未来风险更高？ & 分数、概率、排序和误差评估。\\
因果推断 & 某个行动是否改变结果？ & 处理效应和反事实判断。\\
优化决策 & 在约束下应采取哪些行动？ & 资源分配和策略方案。\\
\bottomrule
\end{tabular}
\end{table}

分析报告应明确模型输出属于风险排序、概率估计、行动效果估计还是资源配置建议，并按相应任务解释结果。

\subsection{标签窗口、特征时间点与样本粒度}
\subjectindex{biaoqianchuangkou}{标签窗口}
\subjectindex{tezhengshidian}{特征时点}

预测问题首先要定义标签。客户流失可以指“未来 30 天取消订阅”，也可以指“未来 90 天没有任何活跃行为”。交付延迟可以按承诺日期判断，也可以按客户实际收货判断。标签定义不同，模型目标和管理行动就不同。

\begin{definition}[标签窗口与特征时间点]
标签窗口是用于观察结果 \(Y\) 的未来时间区间；特征时间点是生成特征 \(X\) 时允许使用信息的截止时点。预测任务要求所有特征都在特征时间点之前可获得，标签则在之后被观察。
\end{definition}

例如，在 2026 年 3 月 1 日预测客户未来 30 天是否流失，则特征只能使用 2026 年 3 月 1 日之前的信息。若使用 3 月 20 日的客服工单预测 3 月内流失，就发生了数据泄漏。只有明确这一时间边界，离线评估才可能反映模型上线后的预测条件。

样本粒度也必须明确。一个客户可以在不同月份形成多条月度记录；以订单为单位时，一张订单对应一条记录；以门店日为单位时，每个门店每天对应一条记录。若训练集和测试集中出现同一客户的高度相似记录，评估结果可能过于乐观。

\section{从线性模型到集成学习}
\subsection{损失函数与经验风险}
\subjectindex{jingyanfengxian}{经验风险}

损失函数规定模型训练的评价目标。若真实结果为 \(y\)，预测为 \(\hat y\)，损失函数 \(\ell(y,\hat y)\) 衡量预测错误的代价。

回归任务常用均方误差：
\[
  \ell(y,\hat y)=(y-\hat y)^2.
\]
分类任务若输出概率 \(\hat p\)，常用二元交叉熵损失：
\[
  \ell(y,\hat p)
  =
  -y\log(\hat p)-(1-y)\log(1-\hat p).
\]
交叉熵会对错误且过度自信的概率预测施加较大惩罚。例如真实 \(y=1\)，若模型给出 \(\hat p=0.01\)，交叉熵损失较大。

统计学习常把模型训练写成经验风险最小化：
\[
  \hat f
  =
  \argmin_{f\in\mathcal F}
  \left\{
  \frac{1}{n}\sum_{i=1}^{n}\ell(y_i,f(x_i))
  +
  \lambda\Omega(f)\right\}.
\]
其中 \(\mathcal F\) 是候选模型类，\(\Omega(f)\) 是复杂度惩罚，\(\lambda\) 控制拟合训练数据和限制模型复杂度之间的权衡。

\begin{modernbox}[管理解释]
损失函数反映预测错误的相对代价。若漏报高风险客户的成本远高于误报，模型评估应报告相应的漏报与误报成本；若留存优惠成本较高，模型阈值就应与预期收益和行动容量共同决定。
\end{modernbox}

\subsection{损失函数与总体风险最小化}
训练数据为目标任务提供有限样本。设希望预测的目标分布为 \(P\)，预测函数为 \(f\)，
其总体风险为 \(R(f)=\E_P[\ell(Y,f(X))]\)。经验风险用样本平均近似这个未知期望；
正则化限制候选函数复杂度。损失函数定义评价标准，函数类限定可表达的关系，
样本数量和分布影响这些关系的估计精度。

以平方损失为例，固定 \(X=x\)，令 \(m(x)=\E[Y\mid X=x]\)，假设二阶矩有限。
对任何数 \(c\)，将 \(Y-c\) 写为 \((Y-m)+(m-c)\) 并展开，交叉项的条件期望为零，
因此
\begin{equation}
\E[(Y-c)^2\mid X=x]
=\operatorname{Var}(Y\mid X=x)+(m(x)-c)^2.
\label{eq:conditional-square-loss}
\end{equation}
第一项不随预测值改变，第二项在 \(c=m(x)\) 时最小。因此，条件均值是平方损失下
的最优预测。更换损失函数，最优预测目标也相应改变。

二元结果同理。若真实条件概率为 \(p\in(0,1)\)，模型报告 \(q\in(0,1)\)，期望对数损失为
\(-p\log q-(1-p)\log(1-q)\)。其对 \(q\) 的导数为
\((q-p)/(q(1-q))\)，在 \(q=p\) 时达到最小。因此，报告真实概率使期望对数损失最小；
仅判断概率位于分类阈值的哪一侧，不能充分评价概率预测的质量。一个模型把0.51和0.99都预测为正类，分类结果
可能一样，概率质量却可能很不相同。当真实概率为0或1时，下确界位于相应边界；
允许报告边界概率时，需按极限解释对数项，不能在开区间内声称存在该边界最小值。

有限样本学习还涉及函数类的逼近误差、参数估计的抽样误差与数值优化误差。
部署环境变化会进一步改变目标分布。区分这些来源，有助于决定是否补充数据、
改善特征、调整函数类或重新估计模型\cite{hastie2009,james2021}。

\subsection{逻辑回归与概率损失的推导}
\subjectindex{luojihuigui}{逻辑回归}
\subjectindex{jiaochashang}{交叉熵}
对二元结果 \(Y\in\{0,1\}\)，令线性分数 \(z=\beta_0+x^\top\beta\)，用 sigmoid 映射到概率 \(p=1/(1+e^{-z})\)。Bernoulli（伯努利）似然\footnote{伯努利分布描述只有两种结果的单次观测，例如流失与未流失；其参数是目标事件发生的概率。似然将已观察结果视为固定，用于比较不同参数值对该结果的支持程度。}为 \(p^y(1-p)^{1-y}\)，取负对数便得到前述交叉熵。关于分数的导数为
\[
\frac{\partial\ell}{\partial z}
=\left(-\frac yp+\frac{1-y}{1-p}\right)p(1-p)=p-y.
\]
当实际发生但预测概率很低时，梯度推动分数上升；反之下降。加入 \(L_2\) 惩罚可限制不稳定的大系数。标准化参数只能从训练数据估计，系数的业务解释还依赖变量编码和共线性\footnote{共线性指多个特征之间存在较强的线性关系，可能使单个系数对样本变化敏感，从而影响系数的稳定解释。}。

平方损失最优预测是条件均值，绝对损失最优预测是条件中位数，分位数损失
\[
\ell_q(y,\hat y)=(q-\mathbf1\{y<\hat y\})(y-\hat y)
\]
对应条件 \(q\) 分位数。它们回答不同问题。若交付承诺关心 \(90\%\) 的订单能否按期完成，直接预测均值不足以替代尾部分位数。

\subsection{随机森林的集成方差}
\subjectindex{suijisenlin}{随机森林}
\subjectindex{jichengxuexi}{集成学习}
设 \(B\) 个基学习器在某个固定输入处的预测误差方差均为 \(\sigma^2\)，两两相关系数为 \(\rho\)。平均后的误差方差为
\[
\operatorname{Var}\!\left(\frac1B\sum_{b=1}^B Z_b\right)
=\sigma^2\left(\rho+\frac{1-\rho}{B}\right).
\]
推导来自 \(B\) 个方差项与 \(B(B-1)\) 个协方差项。在等方差、等相关且 \(\rho\ge0\) 的设定下，增加学习器数量使方差趋于 \(\rho\sigma^2\)。随机森林通过自助抽样\footnote{自助抽样（bootstrap sampling）从原样本中有放回地抽取观测，形成与原样本规模相同或另行指定规模的新样本；随机森林据此训练不同的树。}和节点分裂时的随机特征选择降低树之间的相关性。其总体预测误差还包含偏差，需要在独立数据上评价。

\subsection{梯度提升的迭代更新}
\subjectindex{tidutisheng}{梯度提升}
梯度提升从常数预测 \(F_0\) 开始。在第 \(m\) 步，为每个训练样本计算负梯度
\[
r_{im}=-\left.\frac{\partial\ell(y_i,F)}{\partial F}
\right|_{F=F_{m-1}(x_i)},
\]
用浅树 \(h_m\) 拟合这些目标，再更新
\[
F_m(x)=F_{m-1}(x)+\eta h_m(x),\qquad 0<\eta\le1.
\]
对半平方损失，负梯度就是残差；对 logistic 损失，则是 \(y_i-p_i\)。学习率、树深和早停\footnote{早停是依据验证误差等预设标准，在训练误差继续下降而验证表现不再改善时停止训练，以限制过拟合。}共同限制复杂度。完整实现还需处理分裂增益、正则化、缺失值和计算优化。梯度提升常作为表格数据预测的比较模型，与逻辑回归及基础模型在相同数据划分下评价。

\section{泛化评估与时间切分}
\subsection{训练集、验证集、测试集与数据泄漏}
\subjectindex{shujuxielou}{数据泄漏}

泛化评价使用未参与模型拟合与选择的数据。数据通常按用途划分为以下三部分。
\begin{enumerate}
  \item 训练集 \(\mathcal D_{\mathrm{train}}\)：用于估计模型参数；
  \item 验证集 \(\mathcal D_{\mathrm{valid}}\)：用于选择模型、调参和阈值；
  \item 测试集 \(\mathcal D_{\mathrm{test}}\)：用于最终评估泛化表现。
\end{enumerate}

本章算例采用预先给定的模型设置，使用训练集估计参数、测试集评价结果。需要选择模型或调整超参数时，可在训练资料内部使用验证集或交叉验证。时间序列任务采用早期数据训练、后期数据验证或测试。

\begin{definition}[数据泄漏]
数据泄漏是指模型训练或评估过程中使用了在真实预测时不可获得的信息，从而使离线评估高估模型的未来表现。
\end{definition}

常见数据泄漏包括：用流失之后的客服记录预测流失，用发货之后才知道的实际延迟预测发货风险，用全量数据标准化训练集和测试集，或者在测试集上反复调参。数据泄漏会高估离线表现，模型部署后可能因相应信息不可获得而失效。

\subsection{过拟合、欠拟合与偏差-方差权衡}
\subjectindex{guonihe}{过拟合}
\subjectindex{pianchafangcha}{偏差--方差分解}
\subjectindex{jiaochayanzheng}{交叉验证}

模型太简单时可能欠拟合，太复杂时可能过拟合。对偏差--方差分解，固定输入 \(x\)，设新结果 \(Y=f(x)+\epsilon\)，其中 \(\E[\epsilon\mid x]=0\)、方差为 \(\sigma^2(x)\)，且新噪声与随机训练集 \(\mathcal D\) 独立。则
\[
  \E_{\mathcal D,\epsilon}[(Y-\hat f_{\mathcal D}(x))^2\mid x]
  =
  (\E_{\mathcal D}[\hat f_{\mathcal D}(x)]-f(x))^2
  +
  \operatorname{Var}_{\mathcal D}[\hat f_{\mathcal D}(x)]
  +
  \sigma^2(x).
\]
这个分解区分了系统性偏离、训练样本波动和结果本身的噪声。模型约束较强时可能具有较大偏差，灵活模型在数据不足时可能具有较大方差；但二者并不随模型复杂度必然单调变化，正则化、训练方式和数据规模都会影响结果。模型选择的依据应是与部署条件相符的新样本表现，而非仅凭复杂程度判断优劣。

\begin{definition}[交叉验证]
交叉验证将训练数据分成 \(K\) 份，每次用其中 \(K-1\) 份训练、剩余 1 份验证，循环 \(K\) 次后取平均验证误差。它用于估计模型在未见样本上的表现，并辅助模型选择和调参。
\end{definition}

交叉验证用于模型选择，保留的测试集用于最终评价。若样本具有时间顺序、客户重复观测或门店聚类结构，划分方式应与部署目标一致。若目标是预测全新客户，应按客户分组隔离；若目标是预测已知客户未来一周，则允许使用其过去信息，但必须按时间隔离并处理标签窗口重叠。两种划分分别评价跨客户泛化和跨时间泛化。

\section{概率、排序与行动成本}
\subsection{分类模型的阈值指标}

二元分类模型通常输出概率 \(\hat p_i\)。若选择阈值 \(c\)，则
\[
  \hat y_i=
  \begin{cases}
    1,& \hat p_i\ge c,\\
    0,& \hat p_i<c.
  \end{cases}
\]
于是可以得到混淆矩阵（confusion matrix）\footnote{TP、FP、TN、FN 分别表示真阳性、假阳性、真阴性和假阴性的样本数。阳性表示目标事件发生；假阳性是误报，假阴性是漏报。}：

\begin{table}[htbp]
\centering
\caption{二元分类的混淆矩阵；正类表示目标事件发生}
\begin{tabular}{lcc}
\toprule
\rowcolor{BDMTableHead}
 & 预测为正 & 预测为负\\
\midrule
实际为正 & TP & FN\\
实际为负 & FP & TN\\
\bottomrule
\end{tabular}
\end{table}

常用指标为
\[
  \operatorname{Accuracy}=\frac{TP+TN}{TP+FP+TN+FN},
\]
\[
  \operatorname{Precision}=\frac{TP}{TP+FP},
  \qquad
  \operatorname{Recall}=\frac{TP}{TP+FN},
\]
\[
  F1=
  \frac{2\cdot \operatorname{Precision}\cdot \operatorname{Recall}}
       {\operatorname{Precision}+\operatorname{Recall}}.
\]

准确率（accuracy）衡量总体预测正确比例，但类别不平衡时，单独使用该指标容易产生误导。若真实流失客户只有 \(5\%\)，一个永远预测“不流失”的模型也有 \(95\%\) accuracy，却完全无法支持客户留存。精确率（precision）衡量被标记为高风险的客户中实际流失的比例，召回率（recall）衡量实际流失客户中被模型检出的比例。两者体现了误报和漏报之间的管理权衡。

上述比值在相应分母非零时才有定义。若没有预测正例，precision 的分母为零；
若样本没有实际正例，recall 的分母为零。软件可能约定返回0，但报告应明确这是
计算约定而非观察到的比例。当 \(2TP+FP+FN>0\) 时，可直接计算
\(F1=2TP/(2TP+FP+FN)\)，避免在 precision 与 recall 均为0时再次除以零。

\subsection{排序指标、概率损失与校准}
\subjectindex{AUC}{AUC}
\subjectindex{gailvjiaozhun}{概率校准}

风险排序用于确定复核顺序。排序能力常用 ROC 曲线下面积（AUC）评价\footnote{ROC 是 receiver operating characteristic 的缩写，中文通常称受试者工作特征曲线。它记录不同分类阈值下的假阳性率与真阳性率；AUC 是该曲线下的面积。参见 \url{https://scikit-learn.org/stable/modules/model_evaluation.html}。}。例如，客户经理可以先检查高风险对象，再结合预计行动收益决定联系谁。随机独立抽取一个正样本与一个负样本，记其分数为 \(s_+,s_-\)，则
\[
\mathrm{AUC}=\Prob(s_+>s_-)+\tfrac12\Prob(s_+=s_-).
\]
AUC 不依赖某个固定分类阈值；并列分数按半次排序正确计入。风险排序与行动价值排序仍须分别评价。

AUC 衡量排序能力，收益计算则依赖概率的校准程度。若模型把 \(0.60\) 的概率系统性高估为 \(0.90\)，排序可能仍然正确，但预期收益会被夸大。因此，分类模型还需要检查对数损失（log loss，即前述交叉熵损失）和概率校准。

\begin{definition}[概率校准]
若所有被预测为 \(0.7\) 风险的客户中，长期约有 \(70\%\) 真的流失，则称模型在该概率附近是校准的。校准关注预测概率的数值含义，而不仅是排序。
\end{definition}

校准表通常把预测概率分箱，比较每个箱内的平均预测概率和实际发生率。小样本校准表波动较大，解释时应考虑估计不确定性。概率模型用于预算、收益和风险计算前，需要检查校准程度。

\subsection{阈值选择与业务损失}
\subjectindex{xingdongyuzhi}{行动阈值}

预测模型进入管理流程时，需要把概率转化为行动。若对客户 \(i\) 的留存行动成本为 \(C_i\)，避免流失的价值为 \(M_i\)，并额外假定行动使基线流失风险按相同比例 \(r\) 降低，则一个简化的预期净价值为
\[
  \operatorname{ENV}_i
  =
  \hat p_i \cdot M_i \cdot r - C_i.
\]
当 \(\operatorname{ENV}_i>0\) 时，可以将客户纳入行动候选清单。预测阈值还取决于行动成本和预期收益。高消费客户即使风险略低，也可能值得复核；低消费客户即使风险很高，也未必适合发放高成本优惠。

这里 \(r\) 是相对风险降低比例，不能同时被解释为绝对概率差。若已估计绝对因果降幅 \(\tau_i=\Prob(Y(0)=1\mid X_i)-\Prob(Y(1)=1\mid X_i)\)，正确表达为 \(M_i\tau_i-C_i\)，无需再乘基线概率。真实效果可能随客户变化，第~\ref{ch:week08-causality}~章将估计这一差异。

若行动仅表示分类处置而不改变对象的真实状态，则考虑另一种损失结构：假阳性成本为 \(c_{\rm FP}\)，假阴性成本为 \(c_{\rm FN}\)，正确分类成本为零。比较正类行动的 \(c_{\rm FP}(1-p)\) 与负类行动的 \(c_{\rm FN}p\)，阈值为
\[
p\ge\frac{c_{\rm FP}}{c_{\rm FP}+c_{\rm FN}}.
\]
该式适用于真实状态固定的分类损失；留存干预会改变结果，须使用相应的干预效果模型。

\section{预测区间与覆盖率}

\subsection{分割共形预测}
\subjectindex{gongxingyuce}{共形预测}
预测未来批次的生产周期 \(\hat f(x)\) 时，除了点预测，还需要描述未来结果的不确定范围。共形预测提供一种将模型误差转成预测集合的方法\cite{conformal07}。先将训练、校准与测试数据分开，保持校准样本不参与拟合。再在 \(n\) 个校准样本上计算绝对残差 \(R_i=\abs{Y_i-\hat f(X_i)}\)。然后令
\[
k=\lceil(n+1)(1-\alpha)\rceil,\qquad
q_\alpha=R_{(k)},
\]
其中 \(R_{(k)}\) 是第 \(k\) 个顺序统计量；若 \(k>n\)，约定 \(q_\alpha=+\infty\)。最后输出
\[
C(x)=[\hat f(x)-q_\alpha,\hat f(x)+q_\alpha].
\]
给定训练数据，若校准样本与新样本可交换，则
\[
\Prob\{Y_{\rm new}\in C(X_{\rm new})\}\ge1-\alpha.
\]
可交换是指联合分布在样本次序的任意置换下保持不变。给定训练数据且残差无并列时，
新残差在 \(n+1\) 个残差中的秩服从离散均匀分布。因此，当 \(k\le n\) 时覆盖概率为
\(k/(n+1)\ge1-\alpha\)；存在并列时，包含端点的区间可得到保守覆盖。
上述概率同时对校准样本和未来样本取平均，称为边际覆盖。按设备或产品分别要求的
条件覆盖，需要另行分析。

\begin{example}[有限样本分位数]
取 \(n=9,\alpha=0.2\)，排序残差为
\(1,1,2,2,3,3,4,5,8\) 小时，则 \(k=8\)，区间半径为 5 小时。若预测周期为 20 小时，输出为 \([15,25]\)。这里采用第8个残差，\(n+1\) 修正来自新观测与校准观测的联合秩分布。
\end{example}

\begin{figure}[htbp]
\centering
\begin{tikzpicture}[node distance=0.65cm]
\node[bdmnode] (a) {训练集\\拟合模型};
\node[bdmnode,right=of a] (b) {校准集\\残差分位数};
\node[bdmnode,right=of b] (c) {未来对象\\点预测与区间};
\node[bdmnode,below=of b] (d) {滚动回测\\覆盖与区间宽度};
\draw[bdmarrow] (a)--(b);
\draw[bdmarrow] (b)--(c);
\draw[bdmarrow] (c)|-(d);
\end{tikzpicture}
\caption{校准集与训练集承担不同任务；未来反馈用于监控适用条件。}
\end{figure}

生产时间序列可能具有相关性和分布漂移。应用共形方法时，应检验交换性与部署情形的对应关系，并报告滚动时间划分下的覆盖率、分组覆盖率和区间宽度。加权或自适应方法采用各自的分布与更新条件，区间评价应同时考虑覆盖率和精度。

\section{模型解释与特征贡献}

模型解释研究输入特征如何影响模型输出，其结果用于检查计算行为、发现数据问题和形成后续研究假设。

客户风险模型的输出解释有两类用途：一类是调试模型，检查它是否依赖可疑特征；另一类是支持行动，让客户经理知道应优先处理哪些问题。逻辑回归的系数可以解释为标准化特征对对数优势（log-odds）\footnote{若事件概率为 \(p\)，优势（odds）为 \(p/(1-p)\)，对数优势为其自然对数。逻辑回归假定对数优势是输入特征的线性函数。}的线性贡献，这种贡献描述模型内部的条件关系，干预效应另需因果识别。

\begin{table}[htbp]
\centering
\caption{特征解释中的常见问题}
\begin{tabular}{P{0.24\linewidth}P{0.34\linewidth}P{0.30\linewidth}}
\toprule
\rowcolor{BDMTableHead}
问题 & 合理用途 & 风险\\
\midrule

系数方向 & 检查模型是否符合基本业务直觉。 & 把相关关系解释为因果作用。\\
单客户贡献 & 帮助客户经理理解高风险来源。 & 忽略特征之间的相关性。\\
重要性排序 & 发现模型主要依赖哪些信息。 & 用不稳定小样本重要性做组织决策。\\
异常解释 & 发现数据录入或分布漂移。 & 把数据错误当作真实风险。\\
\bottomrule
\end{tabular}
\end{table}

在客户流失场景中，客服工单多、使用会话少和履约延迟多可以作为行动线索：客户经理可检查服务问题、产品使用问题和交付问题。但这些行动是否能减少流失，需要后续实验或因果评估验证。

\section{前沿模型与企业预测平台}

\subsection{表格与时间序列基础模型}
\subjectindex{TabPFN}{TabPFN}
\subjectindex{Chronos}{Chronos}
\subjectindex{jichumoxing}{基础模型}
2025 年 Nature 的 TabPFN 研究将预训练和上下文学习用于中小规模表格预测；后续 TabPFN-2.5 研究继续扩展这一方向\cite{tabpfn07,tabpfn25}。这类方法将预训练任务中获得的预测能力迁移至新的表格任务。\footnote{预训练是在多个既有任务或大规模数据上先估计模型参数；上下文学习则利用当前输入中的示例适应新任务，而不在该次推理中更新模型参数。}

Chronos-2 将时间序列基础模型扩展到多变量与协变量预测\cite{chronos07}。与传统模型比较时，应固定目标任务、时间划分和损失函数，并报告样本规模、变量类型、推理资源与数据处理方式。

一个可复现的比较顺序是：先固定预测时点和未来窗口；再建立常数或季节性朴素基线；之后比较线性模型、梯度提升与基础模型；最后在相同时间切分上报告预测损失、校准、区间和运行成本。模型选择在训练集与验证集上完成，固定方案后使用保留的测试集报告最终结果。

\subsection{航空备件案例的预测任务拆解}
第~\ref{ch:week05-statistics-dashboard}~章的 Rolls-Royce 公开案例说明发动机与服务网络数据可支持库存规划。相应的预测任务包括故障概率、剩余使用寿命、备件需求量和维修周转时间。它们的标签、删失机制\footnote{删失指结果发生时间未被完整观察，例如研究结束时设备仍未故障，只能知道寿命超过已观察时长。删失记录不能直接当作从未发生故障。}和损失函数不同，备件采购还需结合故障对应的零件需求、维修周期及服务水平确定。

可将样本定义为“发动机在某日的预测快照”，特征来自当时已知的运行和维护记录，标签来自未来固定窗口。提前退役或转出观测范围会产生删失，应记录实际观察时长。生命周期资料较完整时可研究剩余寿命；固定窗口内结果完整时可研究该窗口的故障概率。

\subsection{人工智能辅助建模的结果核查}
AI 工具可以生成特征候选、解释报错和辅助实现模型。模型评价记录应包括数据切分、特征时间证据、基线、训练配置及独立测试结果。智能体自动尝试多个模型时，应保留完整搜索记录，避免只报告表现最好的一次。模型搜索和选择应遵守验证集与测试集的使用边界。

\subsection{模型服务的容量与部署选择}
\subjectindex{KVhuancun}{KV缓存}
企业对语言模型的选取除回答质量外，还取决于输入长度、输出长度、并发和数据
位置。对于具有 \(P\) 个参数、以 \(b_w\) 位存储权重的模型，\(Pb_w/8\) 仅是
权重所需字节的近似下界。若解码器有 \(L\) 层，每层有 \(H_{kv}\) 个键值头，
每头维度为 \(d\)，当前保持 \(n\) 个序列、每序列 \(S\) 个 token\footnote{Token 是分词器处理文本的基本单位，可能对应一个字、一个词或词的一部分，不等同于固定字数。键值缓存（KV cache）保存注意力计算中已生成位置的键和值，以减少后续解码的重复计算。}，键和值
各以 \(b_{kv}\) 字节存储，则未计管理开销的缓存约为
\[
  M_{kv}=2LH_{kv}dSnb_{kv}.
\]
公式中的头数须采用键值头数，不能在分组查询注意力模型中直接换成查询头数。
运行仍有激活、工作空间和碎片等开销；量化权重不会自动按同一比例降低键值
缓存。若 \(L=32,H_{kv}=8,d=128,S=4096,n=8,b_{kv}=2\)，该项为
\(2^{32}\) 字节，即4 GiB。这个数值说明并发与上下文长度会改变部署能力，
上述参数为教学设定\cite{aiinfra2026}。

比较本地模型与云端模型时，先在同一套任务、数据权限和结果判定规则下测量
正确率、平均及尾部时延、每任务费用和运维工作量。令方案 \(j\) 的错误损失为
\(R_j\)、完成时间为 \(T_j\)、单位成本为 \(C_j\)，管理者可在满足
\(R_j\le R_{\max}\)、\(T_j\le T_{\max}\) 及数据位置要求的方案中最小化
\(C_j\)。\(R_j\) 必须按实际业务错误加权；检索型问答、统计预测和工具
执行任务使用不同的判定规则。第~\ref{ch:week09-optimization}~章将把容量与
成本放入资源配置模型。

\section{计算实例与结果解释}

本章样例数据位于：
\begin{center}
\path{data/sample/customer_churn.csv}
\end{center}
数据粒度为一行一个客户观测。字段包括客户编号、训练集与测试集划分、客户分层、客户关系时长、月消费、90 天客服工单数、折扣率、30 天使用会话数、90 天履约延迟次数和是否流失。

配套代码位于：
\begin{center}
\path{src/bdm_decision/cases/week07_prediction.py}
\end{center}
模型工具位于：
\begin{center}
\path{src/bdm_decision/models/churn_model.py}
\qquad
\path{src/bdm_decision/models/evaluate.py}
\end{center}

代码使用一个标准化逻辑回归模型：
\[
  \hat p_i
  =
  \frac{1}{1+\exp[-(\beta_0+\beta^\top z_i)]},
\]
其中 \(z_i\) 是标准化后的特征。训练后的主要系数为：
\begin{table}[htbp]
\centering
\caption{教学流失模型的主要系数；连续特征按训练集标准化}
\begin{tabular}{lr}
\toprule
\rowcolor{BDMTableHead}
特征 & 系数\\
\midrule
客户关系时长 & -2.2397\\
90 天客服工单数 & 1.1511\\
30 天使用会话数 & -1.0217\\
折扣率 & 0.8472\\
90 天履约延迟次数 & 0.6323\\
月消费 & 0.3712\\
\bottomrule
\end{tabular}
\end{table}
在该模型中，客户关系较短、客服工单较多、使用活跃度较低、折扣较高和履约延迟较多均对应较高的预测流失风险。这些系数刻画给定其他输入时的模型输出变化，干预效应需另行识别。

在测试集上，阈值 \(0.5\) 下得到：
\begin{itemize}
  \item 混淆矩阵：TP=5，FP=1，TN=4，FN=0；
  \item accuracy 为 0.9000；
  \item precision 为 0.8333；
  \item recall 为 1.0000；
  \item F1 为 0.9091；
  \item ROC AUC 为 1.0000；
  \item log loss 为 0.1397。
\end{itemize}

最高风险客户为 C030、C028、C021、C025 和 C023。若进一步考虑月消费、三个月风险毛利、留存行动效果和优惠成本，预期净价值排序最高的客户是 C028、C025、C023 和 C021；C030 虽然风险最高，但月消费较低，按当前成本假设不建议发放同样的优惠。风险排序与预期净价值排序在该算例中不同。

\section{模型运行与风险管理}
\subjectindex{fenbupianyi}{分布偏移}

上线监控应同时检查输入变量与标签生成策略的变化。假设历史上
高风险订单通常会被人工加急，那么训练标签反映的是“历史人工策略作用之后”的
延期状态。模型一旦改变加急名单，未来标签的生成过程
也随之改变，原有条件关系可能不再保持。

在满足一致性、给定特征无混杂等条件的简化情形中，令历史加急概率为 \(e(x)\)，
两种行动的条件结果均值为 \(m_1(x)\) 与 \(m_0(x)\)，则历史观察均值满足
\[
\E[Y\mid X=x]=e(x)m_1(x)+(1-e(x))m_0(x).
\]
这个混合关系说明，即使两个结果模型不变，改变行动策略 \(e(x)\) 也会改变观察到的
平均风险。若无混杂不成立，两个条件分量还不能解释为潜在结果均值，问题更加复杂。

因此，模型发布记录应保存当时的行动政策与人工干预日志，而不只保存训练数据和
模型文件。评价也应区分风险预测是否失准、干预是否奏效，以及处理分配是否使更难的
对象集中在未处理组。行动效果的比较需要识别相应的反事实关系，识别条件见
第~\ref{ch:week08-causality}~章。

模型部署后的适用性取决于当前业务环境与训练期环境的关系。客户行为、产品价格、营销活动、竞争环境和数据采集方式都会变化。模型上线后至少应监控以下事项：
\begin{enumerate}
  \item 数据漂移：特征分布是否偏离训练期；
  \item 标签漂移：流失率、延期率或需求水平是否发生结构变化；
  \item 性能漂移：precision、recall、log loss 和校准是否变差；
  \item 决策策略变化：模型推荐是否改变了后续数据生成过程；
  \item 公平性与合规：不同客户分层、地区或渠道是否受到不合理影响；
  \item 结果跟踪：被干预对象的结果是否被记录，用于评估行动效果。
\end{enumerate}

运行监测还需考虑业务处理能力。若客户经理只能联系 20 个客户，每日候选名单需要根据收益、处理容量和时限筛选。上线系统必须同时说明行动容量、责任人、处理时限和复盘机制。

\subsection{模型风险管理}

预测模型进入企业流程前，应至少检查以下风险：
\begin{enumerate}
  \item 估计精度：测试样本量是否足以区分候选模型的预测误差；
  \item 样本选择偏差：训练客户是否来自特定渠道、地区、时间段或分层；
  \item 数据泄漏：特征是否包含预测时不可获得的信息；
  \item 标签定义：流失窗口、暂停、降级、退款和自然到期是否被统一处理；
  \item 概率校准：预测概率能否用于预算和收益计算；
  \item 阈值治理：阈值是否由成本、容量、风险偏好和公平性共同决定；
  \item 结果跟踪：被干预客户的后续结果是否记录，用于评估行动效果；
  \item 模型漂移：客户行为、产品策略和市场环境变化后，模型是否仍然有效。
\end{enumerate}

预测模型报告应包括数据切分、特征字典、模型表达、训练设置、测试指标、校准检查、阈值说明、行动建议和风险备注，作为模型审计与运行维护的依据。

\section*{本章小结}
\addcontentsline{toc}{section}{本章小结}

监督学习利用带标签历史数据估计从状态 \(X\) 到未来结果 \(Y\) 的预测函数；损失函数和经验风险定义了模型训练目标；训练集、验证集、测试集和交叉验证用于控制泛化风险；分类模型需要同时报告阈值指标、排序指标、概率损失和校准检查。预测概率用于描述风险。行动还需要考虑成本、收益、容量、干预效果、公平性和治理机制。

逻辑回归由二元似然得到概率损失，集成方法结合多个学习器改善预测，共形方法在交换性条件下构造具有边际覆盖性质的预测集合。运行评价还须跟踪目标分布与决策策略的变化；语言模型的部署另需核算权重、键值缓存、并发和输出质量。行动对结果的影响需依据潜在结果框架研究，见第8章。

\chapterreading{统计学习的基础推导可参考 Hastie 等、James 等及 Bishop
\cite{hastie2009,james2021,bishop2006}，概率建模背景见 Murphy\cite{murphy2012}。
Breiman 关于两种建模文化的讨论有助于理解预测目标与机制解释的差异
\cite{breiman2001}，ROC 评价见 Fawcett\cite{fawcett2006}。共形预测的系统入门见
Angelopoulos 与 Bates\cite{conformal07}。表格基础模型的有限样本表现仍需在明确标签定义、时间切分与部署条件后评价。语言模型的负载与资源估算见\cite{aiinfra2026}。}

\section*{思考与习题}
\addcontentsline{toc}{section}{思考与习题}

\begin{exercise}[集成与区间]
取基模型方差为 1、相关系数为 0.3，比较 1、10、100 个模型平均后的方差。再重算本章共形区间，并说明在产品换代导致分布突变后为什么不能照搬覆盖保证。
\end{exercise}
\begin{exercise}[风险与效果]
客户甲的流失风险为 \(0.8\)、可降低的绝对概率为 \(0.02\)，客户乙分别为 \(0.4\) 与 \(0.15\)。两者避免流失价值均为 1000、行动成本均为 50，计算净价值并解释为什么风险排序和行动排序相反。
\end{exercise}

\begin{exercise}
围绕 \path{data/sample/customer_churn.csv} 和 \path{src/bdm_decision/cases/week07_prediction.py} 完成下列任务：
\begin{enumerate}
  \item 说明每个特征在客户流失预测中的业务含义，并指出它是否可能存在数据泄漏；
  \item 运行第7章脚本，记录测试集 confusion matrix、accuracy、precision、recall、F1、ROC AUC 和 log loss；
  \item 将阈值从 0.50 改为 0.70，比较误报和漏报如何变化；
  \item 选择 3 个最高风险客户，说明他们是否也一定是最值得留存行动的客户；
  \item 为该模型设计一张上线监控表，至少包含数据漂移、性能漂移和行动容量；
  \item 写一段不超过 300 字的管理建议，必须同时提及模型指标、业务成本和模型风险。
\end{enumerate}
\end{exercise}
